%--IPFW Systems Engineering Course Template--%
%--Maintained by David S. Cochran and Joesph J. Smith
% Email questions/concerns to smitjj09@students.ipfw.edu
\documentclass{ipfwpaper}
\svnid{$Id$}


%%%%%------Instructions------%%%%%%%%

% Step 1: Add Cover Page Details
	%--- Add your names and emails
    \author{Author 1\formatemail{Email 1} Author 2\formatemail{Email 2} Author 3\formatemail{Email 3}}
    
% Step 2: Add Course Number (i.e., SE540, SE550, etc..)
    \course{SE5XX}
    
% Step 3: Add your title
	\title{Systems Engineering Project Template} 
    
% Step 4: Format ACRONYMS using Latex, see the Overleaf primer and the acronyms.tex file for instructions.

	% Use the acronyms.tex file (included under 		Project 	
	% files) to format acronyms.
    
% Step 5: Fill in the remaining sections with content
	%----This document contains content for various SE Course
	%--- Uncomment the Section related to your course 
    %    (Highlight section, and hit "Ctrl" and "/" keys simultaneously)
	%--- Refer to the Overleaf Primer for formatting help


%----Acronyms-----% 
\loadglsentries{acronyms} %Formats Acronyms
\makeglossaries %Makes Glossary
%-----------------%

%%%%%%%%%%%%%%%%%%%%%%%%%%%%%%%%%%%%%%%%%%%%%%%%
%-------------Beginning of Document------------%
%%%%%%%%%%%%%%%%%%%%%%%%%%%%%%%%%%%%%%%%%%%%%%%%
\begin{document} 
\maketitle %Makes the cover page


%---Revision Tracker---%
\startTable
\AddRevision{Number}{Date}{Add explanation}
\AddRevision{}{}{}
\stopTable

%-----Table of Contents-----%
\tableofcontents
\listoftables
\listoffigures
\printglossary[title=List of Acronyms ,type=\acronymtype]
%--end:Table of Contents--%




%%%%%%%%   SE 52000 - Advanced Manufacturing and Processes %%%%%%%
 \section{Abstract}  %Short explanation of project
 \section{Introduction} 
 
\section{Literature Review}
This section should includes information related to: 
\begin{itemize}
    \item What others are doing in this field
    \item ERP Software background and uses
    \item Applicable DoD Regulations
\end{itemize}

    \section{Problem Statement}
    State what the project is about and the difficulties that may be involved.
	\subsection{Customer(s)}
	Customers can be both internal and external to the system, and it is important to meet the needs of both types. Refer to the System Design Language Primer for more information. 
    	\subsubsection{Types of Customers}
        \subsubsection{Customer Needs}
        \subsection{Use Case(s) or Scenario(s)} 
         Explain the function of the product and enterprise from a viewpoint of a customer. What functions must the product and enterprise achieve in order to satisfy the needs of the various customers? Defining use cases comes prior to establishing the system boundary because it may help understand the scope of the project. 
    	\subsection{System Boundary}
	   The system boundary is usually expressed with a System Boundary Diagram. Please see the System Design Language Primer for more information.
    \subsection{Interfaces and Interface Risks}
    
    
\section{Product Design Description}
Describe your team's product and be sure to explain the differences between the commercial product and the Defense-related product, if applicable.

\section{Manufacturing System Design}
	\subsection{Translation of Customer Needs (CNs) to Candidate Functional  Requirements (FRs)}
    \subsection{Manufacturing System Design Decomposition/Map}
    \subsection{Value Streams}
    	\subsubsection{Description of Value Streams/Major Process Steps}
    	Please define the FRs and PSs of your manufacturing system.
    	\subsubsection{Part Family Definition and Overall Linked-Cell Layout for the Value Stream}
    	\subsubsection{Value Stream Map}
    	The Value Stream Map should show the material and information flow relationship as well as how the cell cycle time is balanced to takt time.
\section{Company Structure}
    \subsection{Systems Engineering Management Plan (SEMP)}
    \subsection{DCAA Audit Compliance}
         Please include your answers to the SF1408 Form in this section.
    \subsection{ERP Software}
        \subsubsection{Use Case Description}
        \subsubsection{ERP Software Specification Requirements}
    \subsection{Work Breakdown Structure (WBS)}
    \subsection{Integrated Master Schedule}
    \subsection{Chart of Accounts}
    \subsection{Budget}
    
\section{Pre-Meeting: DCAA Preparation Materials}
    Include any material that your team has created in preparation for the DCAA Meeting 1. This pre-meeting 1 is used for review and to allow your team to get any questions answered prior to the first presentation. 
    
\section{DCAA: Meeting 1}
    Include the content and description for the DCAA Meeting 1.
    
\section{ERP Software: Meeting 2}
    Include the content and description for the ERP Software Meeting 2.
    
\section{Board of Directors: Meeting 3} 
    Include the content and description for the Board of Directors Meeting 3.
    
 \clearpage      
\section{Conclusions}
\section{Lessons Learned}

\section{References}
    \bibliographystyle{IEEEtran}
    \bibliography{your_references.bib}
\appendix





       
% %%%%%%%%   SE 53000 - Systems Engineering Management  %%%%%%%
%  \section{Abstract}  %Short explanation of project
%  \section{Introduction} %Explain topic as though the reader has no knowledge of the subject
 
% \section{Literature Review}%What is the current research on the topic?                                    
% \section{Problem Statement}
% 	\subsection{Customer(s)}
%     	\subsubsection{Types of Customers}
%         \subsubsection{Customer Needs}
% 	\subsection{Use Case(s) or Scenario(s)}                     
% 	\subsection{System Boundary}
%     \subsection{Interfaces and Interface Risks}
% \section{System Design}
% 	\subsection{Translation of Customer Needs (CNs) to Candidate Functional  Requirements (FRs)}
%     \subsection{System Design Map (to 5 levels)}
%     \subsection{Value Streams}
%     	\subsubsection{Description of Value Streams/Major Process Steps}
%     	\subsubsection{Activities, Connection and Flow (ACF) Diagram for 						Each Process/Value Stream}
%     \subsection{Standard Work}
%     	\subsubsection{Job Classification Requiring Standard Work}
%     	\subsubsection{Standard Work for each Job Classification/Activity}
%         \subsubsection{Process for Sustaining Standard Work}
%         \subsubsection{Process for Improving Standard Work} 
%         	%Include TONE for Continuous Improvement
%  \clearpage      
% \section{Conclusions}
% \section{Lessons Learned}
% 	\subsection{Related to Design}
%     	\subsubsection{Use of Collective System Design}
%     	\subsubsection{Use of Lean}
%     \subsection{Related to People's Job Definition}
%     \subsection{Related to Improvement}
	
% \section{References}
%     \bibliographystyle{IEEEtran}
%     \bibliography{your_references.bib}
% \appendix
% 	\section{System Design Map}
  
 
    
% %%%%%%%%%   SE 54000 - System Architecture    %%%%%%%%% 
% \section{Abstract}  %Short explanation of project

% Text for the abstract will be entered here. Formating commands can be found on the Overleaf Primer. Please be sure to use the acronym function to define acronyms. The acronyms.tex file is used to declare new acronyms, but example ones have been included as well. For example, an acronym you may need for this project is Functional Requirement. Instead of typing out the whole acronym, the command created in the acronyms.tex can be used. For example, the acronyms already defined include; \gls{fr}, \gls{ps}, and \gls{csd}. The next time you use the command, only the abbreviation will be shown. For example, \gls{fr} will only be printed.


% \section{Introduction} %Explain topic as though the reader has no knowledge of the subject
% \section{Literature Review}%What is the current research on the topic?   

% \clearpage
% \section{Problem Statement}
% 	\subsection{Customer(s)}
%     	\subsubsection{Types of Customers}
%         \subsubsection{Customer Needs}
% 	\subsection{Use Case(s) or Scenario(s)}                  
%  	\subsection{System Boundary}
%     \subsection{Interface Definition and Interface Risks}
    
% \clearpage
% \section{System Design}
% 	\subsection{Logical Interfaces and Risks}
%  	\subsection{Translation of Customer Needs (CNs) to Functional  Requirements (FRs)}
%      \subsection{Product Design Decomposition Map}
     
%  \clearpage
% \section{Architecture Framework}
 
%  \clearpage
% \section{Detailed Design}
% 	\subsection{Physical Interfaces and Interface Risks}
% 	\subsection{Design Description and Drawings}
%     \subsection{Materials Definition and Selection}
%     \subsection{Applicable Standards}
% 	\subsection{Analysis and Calculations}
%     \subsection{DFA}
    
%  \clearpage
%  \section{Risk Mitigation, FMEAs, Test Plan}
% 	\subsection{Test Plan}
% 	\subsection{Risk Mitigation Plan}
%     \subsection{Design FMEA}
%     \subsection{Process FMEA}
    
% \clearpage
% \section{Design Enhancements}      
% \section{Conclusions}
% \section{Lessons Learned}

% \clearpage
% \section{References}
%     \bibliographystyle{IEEEtran}
%     \bibliography{your_references.bib}
    
% \clearpage
% \appendix
%  \renewcommand{\thesection}{\Alph{section}.}
% \setcounter{section}{0}

% \addcontentsline{toc}{section}{Appendices}

% %-----Appendix A-----%
% \section{Product Design Decomposition Map}

% %-----Appendix B-----%
% \section{...}



% %%%%%%%%   SE 55000 - Advanced Manufacturing and Processes %%%%%%%
%  \section{Abstract}  %Short explanation of project
%  \section{Introduction} %Explain topic as though the reader has no knowledge of the subject
 
% \section{Literature Review}%What is the current research on the topic?                                    
% \section{Problem Statement}
% 	\subsection{Customer(s)}
%     	\subsubsection{Types of Customers}
%         \subsubsection{Customer Needs}
% 	\subsection{Use Case(s) or Scenario(s)}                     
% 	\subsection{System Boundary}
%     \subsection{Interfaces and Interface Risks}
% \section{Product Design}
% \section{Manufacturing System Design}
% 	\subsection{Translation of Customer Needs (CNs) to Candidate Functional  Requirements (FRs)}
%     \subsection{System Design Map (to 5 levels)}
%     \subsection{Value Streams}
%     	\subsubsection{Description of Value Streams/Major Process Steps}
%     	\subsubsection{Value Stream Map}
%     \subsection{Cell Design}
%     \subsection{Standard Work}
%     	\subsubsection{Standard Work for each Job Classification/Activity}
%         \subsubsection{Process for Sustaining Standard Work}
%         \subsubsection{Process for Improving Standard Work} 
%         	%Include TONE for Continuous Improvement
%  \clearpage      
% \section{Conclusions}
% \section{Lessons Learned}
% 	\subsection{Related to Design}
%     	\subsubsection{Use of Collective System Design}
%     	\subsubsection{Use of Lean}
%     \subsection{Related to Improvement}
	
% \section{References}
%     \bibliographystyle{IEEEtran}
%     \bibliography{your_references.bib}
% \appendix
% 	\section{System Design Map}







\end{document}





